\documentclass[11pt]{article}
\usepackage[a4paper,margin=25mm]{geometry}
\usepackage{amsmath,amssymb}
\usepackage[T1]{fontenc}
\usepackage{hyperref}
\setlength{\emergencystretch}{3em}
\title{Rank and crank are not equidistributed:\\a counterexample to Conjecture 7292}
\author{gaochengzhecpu (AI-assisted submission)}
\date{3 October 2026}
\begin{document}
\maketitle

\paragraph{Statement being disproved.}
Conjecture \texttt{00000007292} in The Last Math Competition asserts that
the rank and crank of integer partitions have the same distribution,
and adds a claim about their generating function. We disprove the
equidistribution clause using partitions of $4$. Here distribution is by
the integer value of the statistic, as in the standard definitions below;
the problem specifies no modulus or restricted arithmetic progression.
Disproving this first clause suffices to disprove the stated conjunction.

\paragraph{Definitions.}
A partition of $n$ is a finite sequence
$\lambda=(\lambda_1,\ldots,\lambda_\ell)$ of positive integers in weakly
decreasing order, with sum $n$. For a nonempty partition, its Dyson rank is
\[
 r(\lambda)=\lambda_1-\ell.
\]
Let $o(\lambda)$ count its parts equal to $1$, and let $u(\lambda)$ count
its parts strictly greater than $o(\lambda)$. The Andrews--Garvan crank is
\[
 c(\lambda)=
 \begin{cases}
 \lambda_1,&o(\lambda)=0,\\
 u(\lambda)-o(\lambda),&o(\lambda)>0.
 \end{cases}
\]
These are the standard statistics used by Atkin and Garvan
\cite[Introduction]{AG}. Put
\[
 N(m,n)=\#\{\lambda\vdash n:r(\lambda)=m\},\qquad
 M(m,n)=\#\{\lambda\vdash n:c(\lambda)=m\}.
\]
Equidistribution would imply $N(m,n)=M(m,n)$ for every integer $m$ and
every positive integer $n$.

\paragraph{Counterexample.}
The complete list at $n=4$, with both statistics computed from their
definitions, is
\[
\begin{array}{c|r|r|r|r}
 \lambda & \ell & o(\lambda) & r(\lambda) & c(\lambda)\\\hline
 (4)       &1&0& 3& 4\\
 (3,1)     &2&1& 1& 0\\
 (2,2)     &2&0& 0& 2\\
 (2,1,1)   &3&2&-1&-2\\
 (1,1,1,1) &4&4&-3&-4
\end{array}
\]
For completeness, a partition of $4$ has at most four parts. With one
part it is $(4)$; with two parts the decreasing positive solutions of
$a+b=4$ are $(3,1)$ and $(2,2)$; with three parts only $(2,1,1)$ is
possible; with four parts all parts equal $1$. This proves the list is
exhaustive. Hence
\[
 N(4,4)=0\ne1=M(4,4),
\]
contradicting equidistribution. The same discrepancy can be seen without
the table: every partition of $4$ has largest part at most $4$ and at
least one part, so its rank is at most $3$, whereas $(4)$ has crank $4$.

\paragraph{Formal verification and semantic bridge.}
The self-contained Lean 4.19.0 project imports only \texttt{Std}.
\texttt{IsPartition n parts} is the standard definition: the list is
weakly decreasing, every part is positive, and its sum is $n$.
\texttt{first\_is\_largest} verifies that the first list entry used in
both statistics is indeed the largest part. Both subtractions take place
in \texttt{Int}, so negative ranks and cranks are retained.

The theorem \texttt{partitionsFour\_complete} quantifies over
\emph{arbitrary lists of natural numbers}, proves their length is at most
four, and classifies each possible length. Together with
\texttt{partitionsFour\_exact} and \texttt{partitionsFour\_nodup}, this
shows that counting the listed objects is precisely counting all standard
partitions once, rather than sampling selected partitions.

\texttt{RankCrankEquidistribution} states equality of all statistic-fiber
counts for every complete duplicate-free enumeration of the partitions
of $n$. Thus its counts are ordinary finite cardinalities, independent of
the chosen enumeration. The final theorem
\texttt{conjecture7292\_counterexample} negates this universal assertion
by applying it to the verified enumeration at $n=4$. The supplementary
theorem \texttt{not\_equal\_supports} independently formulates the
contradiction directly with existential quantifiers over all partitions:
crank $4$ occurs and rank $4$ does not.

The project contains no proof placeholders, native decision procedure,
or added axioms. The displayed \texttt{\#print axioms} audits use only
Lean's standard \texttt{propext}, \texttt{Classical.choice}, and
\texttt{Quot.sound}. There is no dependence on a generating-function
identity, and the exceptional convention sometimes used for crank
generating functions at $n=1$ is irrelevant to $n=4$.

\begin{thebibliography}{9}
\bibitem{AG}
A.~O.~L. Atkin and F.~G. Garvan,
\emph{Relations between the ranks and cranks of partitions},
Ramanujan J. \textbf{7} (2003), 343--366.
\href{https://arxiv.org/abs/math/0208050}{arXiv:math/0208050}.
\end{thebibliography}
\end{document}
